\section{Dividing the work between agents, code and people}

The previous sections describe what agents have been asked to do.
This section turns to a question every one of those systems answers, explicitly or not: which parts of the work go to the language model, which to ordinary code, and which stay with a person.
The most useful evidence comes from recent work on AI-assisted research in economics and the social sciences, where the question has been studied directly.

\subsection{From full autonomy to gated pipelines}

At one end of the field are systems built for full autonomy.
The AI Scientist generates research ideas, writes and runs experiments, and drafts and reviews complete machine-learning papers with little human involvement, at a reported cost of under 15 dollars per paper \citep{lu2024ai}.
Coscientist plans and executes chemistry experiments through laboratory automation \citep{boiko2023autonomous}, ChemCrow gives a language model eighteen expert-designed chemistry tools \citep{bran2024augmenting}, and Co-Scientist generates and ranks biomedical hypotheses for experimental testing \citep{gottweis2026accelerating}.
Surveys of these scientific agents emphasise domain tools and validation mechanisms as what separates them from general-purpose assistants \citep{ren2025towards}.

Evaluations temper the picture.
ScienceAgentBench draws 102 tasks from 44 peer-reviewed publications and asks agents to write the analysis code for each; the best agent solved about a third of the tasks \citep{chen2024scienceagentbench}.
Its authors draw a methodological conclusion that this thesis adopts: agents should be assessed on the individual tasks of a scientific workflow before claims are made about automating the workflow end to end.

In economics, \citet{korinek2023generative} catalogues dozens of uses of LLMs across ideation, writing, background research, coding, data analysis and mathematical derivation, and argues that the main near-term gains come from automating many small tasks within the researcher's own workflow.
A later guide describes how economists can build agents for literature, coding and data work \citep{korinek2025ai}.
Both treat the researcher as the one who decides what is worth doing and whether the result holds.

\subsection{Evidence that the division matters}

The clearest evidence on how to divide the work comes from a pair of papers on human-in-the-loop economic research.
HLER is a multi-agent pipeline for empirical economics that audits a dataset, generates hypotheses constrained by what the data can support, runs the econometric analysis, and drafts a paper, with human decisions at fixed points \citep{zhu2026hler}.
\citet{zhu2026human} then test whether its reliability comes from the model or from how the work is organised.
In a pre-specified experiment with 280 complete research runs on four datasets, an unconstrained multi-agent pipeline produced a critical failure in 72\% of runs.
With the same model, the same agents and the same prompts, the failure rate fell to 16\% when three rules were imposed: the LLM reasons but does not do the data work, data preparation and estimation are done by deterministic code, and a person decides at three gates (the choice of research question, the identification strategy and the decision to publish).
An ablation suggests that deterministic computation and human gates each reduce failures on their own, and do most when combined.
The most common failures in the unconstrained pipeline were hallucinated references, inconsistent interpretation, infeasible research questions and weak identification, not errors in executing code.
The authors note the limits: one model, four datasets, and an ablation with only twenty runs per condition.
Their conclusion is that the question is not how to remove people from the loop, but where human attention does the most to prevent, expose and contain failure.

The same authors address a harder case in pAI-Econ-claude, a gated multi-agent architecture for developing economic theory \citep{zhu2026pai}.
They start from the observation that many outputs in economics have no cheap, complete, machine-readable signal of whether they are correct.
A proof sketch, a mechanism or a welfare claim cannot be checked the way code can be tested.
Their gates therefore \emph{diagnose} specific failures and recommend going back a step, but explicitly do not certify that what passes is correct.
Decisions that are costly to reverse, such as the choice of equilibrium concept, are reserved for the researcher.
In five paired comparisons, two blinded evaluators preferred the gated workflow in four, though it used several times the resources of the ungated baseline.

\subsection{What this strand establishes}

This work offers the thesis two things.
The first is a design principle with experimental support: let the language model reason and propose, let deterministic code compute, and put people at the decisions that are hard to undo.
The engineering systems in Sections 2.2 to 2.5 arrive at the same division independently, through harnesses, validators and solvers that produce every number.
The second is a recognition, now explicit in the AI-for-economics literature, that some outputs cannot be scored against a right answer.
Where that is so, automated checks can find specific faults, but deciding that a result is good enough remains a human judgment.
